% Все тексты урока редактируются здесь. Интерактивная модель: \interactive{conditional-geometry}.
\section{Теория}

\subsection{Условие меняет множество, в котором мы считаем}
Представим множество элементарных исходов $\Omega$ квадратом, разделённым на клетки. В каждой клетке записано, сколько элементарных исходов в ней находится. Все элементарные исходы равновероятны. Каждый исход находится ровно в одной клетке.

Горизонтальная полоса задаёт событие $A$, вертикальная — событие-условие $B$. Все исходы в клетках на пересечении полос составляют $A\cap B$.

Если известно, что произошло $B$, исходы вне вертикальной полосы больше не рассматриваются. Нас интересует, какая доля исходов внутри $B$ находится также в $A$.

\begin{definition}
Обозначим через $n(E)$ число элементарных исходов, из которых состоит событие $E$, а через $N=n(\Omega)$ — общее число исходов. Поскольку все исходы равновероятны,
\[
 P(E)=\frac{n(E)}{N}.
\]
Чтобы найти число исходов в полосе, складываем числа во всех её клетках.
\end{definition}

\subsection{Подвигайте события}
На рисунке число в клетке — количество элементарных исходов в ней. Чем темнее клетка, тем больше исходов она содержит. Зелёная горизонтальная полоса — $A$, золотистая вертикальная — $B$, терракотовое пересечение — $A\cap B$.

Перетаскивайте полосы за внутреннюю область или подпись. Потяните за круглые ручки на границах, чтобы изменить число выбранных строк или столбцов. Полосы перемещаются по целым клеткам. Те же действия доступны ползунками; их можно менять стрелками клавиатуры.

\interactive{conditional-geometry}

\subsection{Почему знаменатель теперь другой}
\begin{theorem}[conditional-probability-geometry/theorem/1]\label{theorem-1}
Если $n(B)>0$, то
\[
 P(A\mid B)=\frac{P(A\cap B)}{P(B)}=\frac{n(A\cap B)}{n(B)}.
\]
В знаменателе стоит число исходов внутри полосы-условия, а не число исходов во всём квадрате.
\end{theorem}

\begin{proof}[conditional-probability-geometry/theorem/1]\label{solution-conditional-probability-geometry-theorem-1}
Подставим вероятности через числа равновероятных исходов в определение условной вероятности:
\[
 P(A\mid B)=\frac{n(A\cap B)/N}{n(B)/N}=\frac{n(A\cap B)}{n(B)}.
\]
Общий делитель $N$ сокращается. После условия $B$ новым целым становится множество исходов внутри полосы $B$.
\end{proof}

\subsection{Что меняется при движении полос}
Перемещайте вертикальную полосу при неизменной горизонтальной. Вероятность $P(A)$ останется прежней: ни событие $A$, ни общее число исходов не изменились. Однако в новое условие $B$ могут попасть клетки с другим количеством исходов, поэтому $P(A\mid B)$ может измениться.

Клетки одинаковой площади могут содержать разное число элементарных исходов. Поэтому одинаковая площадь пересечения ещё не означает одинаковую условную вероятность: считать нужно исходы, попавшие в пересечение, и исходы внутри условия.

Равновероятность элементарных исходов не означает равновероятность клеток. Например, клетка с восемью исходами в восемь раз вероятнее клетки с одним исходом. Форма горизонтальной и вертикальной полос сама по себе не гарантирует независимости событий $A$ и $B$.

\subsection{Если условие невозможно}
Если сжать $B$ до нулевой ширины, в нём не останется исходов и $P(B)=0$. Условная вероятность $P(A\mid B)$ тогда не определена: делить на ноль нельзя. Это не вероятность, равная нулю.
